\documentclass[11pt, a4paper]{article}

\usepackage[T1]{fontenc}
\usepackage[utf8]{inputenc}
\usepackage{tgpagella}
\usepackage{tgcursor}
\usepackage{microtype}
\usepackage[spanish, es-noshorthands]{babel}
\usepackage[table, xcdraw]{xcolor}
\usepackage{booktabs}
\usepackage{tabularx}
\usepackage{array}
\usepackage{amsmath, amssymb}
\usepackage[most]{tcolorbox}
\usepackage{tikz}
\usetikzlibrary{shapes.geometric, arrows.meta, positioning, calc}
\usepackage[margin=2cm]{geometry}
\usepackage{fancyhdr}

\pagestyle{fancy}
\fancyhf{}
\setlength{\headheight}{16pt}
\lhead{\textbf{IMT-342}}
\chead{Matriz de Arquitecturas A/B y Diseño Gate 1}
\rhead{Semana 12 - Sesión 2}
\lfoot{Diseño de Prueba Solo Lectura (Sin Movimiento)}
\rfoot{Página \thepage}
\renewcommand{\headrulewidth}{0.4pt}
\renewcommand{\footrulewidth}{0.4pt}

\definecolor{ucbblue}{RGB}{0, 51, 102}

\begin{document}

\begin{center}
    \Large\textbf{\textcolor{ucbblue}{Contraste de Arquitecturas Candidatas (A/B) y Diseño del Ensayo Gate 1}}[cite: 16]
\end{center}

\section{Separación Conceptual de Canales y Visual 3[cite: 16]}
Para sostener una afirmación de comunicación en Gate 1, es obligatorio distinguir entre un comando emitido desde ROS, el transporte de red, el estado generado por el controlador y el movimiento físico[cite: 16]. Ni la visualización en RViz, ni el eco local de un tópico, ni un \texttt{ping} exitoso prueban que se lean estados reales del controlador ABB[cite: 16].

\begin{center}
% Visual 3: Arquitecturas A y B con separación de lectura/escritura
\begin{tikzpicture}[
    box/.style={rectangle, draw=ucbblue, thick, fill=ucbblue!5, text width=3.0cm, align=center, rounded corners, font=\scriptsize, minimum height=1.1cm},
    bridge/.style={rectangle, draw=orange!80!black, thick, fill=yellow!10, text width=3.6cm, align=center, rounded corners, font=\scriptsize, minimum height=1.1cm},
    hw/.style={rectangle, draw=gray!80!black, thick, fill=gray!10, text width=2.8cm, align=center, rounded corners, font=\scriptsize, minimum height=1.1cm},
    auth/.style={rectangle, draw=red!80!black, very thick, fill=red!5, text width=11.5cm, align=center, rounded corners, font=\scriptsize\bfseries, inner sep=2mm}
]
    \node[box] (ros) at (0, 0) {\textbf{PC / Entorno ROS}\\Nodos de Consulta y Registro[cite: 16]};
    \node[bridge] (drv) at (5.5, 0) {\textbf{Candidato A} (Driver Existente)\\o \textbf{Candidato B} (Puente Adaptado)[cite: 16]};
    \node[hw] (ctrl) at (11.0, 0) {\textbf{Red $\leftrightarrow$ Controlador ABB}\\(Robot Físico en Reposo)[cite: 16]};

    % Canal superior: Solicitud / Comando (Bloqueado para movimiento)
    \draw[-Stealth, thick, red!80!black] ([yshift=2.8mm]ros.east) -- node[above, font=\tiny\bfseries]{Solicitud Consulta (NO mover hoy)[cite: 16]} ([yshift=2.8mm]drv.west);
    \draw[-Stealth, thick, red!80!black] ([yshift=2.8mm]drv.east) -- node[above, font=\tiny\bfseries]{Petición Solo Lectura[cite: 16]} ([yshift=2.8mm]ctrl.west);

    % Canal inferior: Estado validado / Identificación
    \draw[-Stealth, very thick, ucbblue, dashed] ([yshift=-2.8mm]ctrl.west) -- node[below, font=\tiny\bfseries]{Estado / ID + Timestamp[cite: 16]} ([yshift=-2.8mm]drv.east);
    \draw[-Stealth, very thick, ucbblue, dashed] ([yshift=-2.8mm]drv.west) -- node[below, font=\tiny\bfseries]{Realimentación Trazable[cite: 16]} ([yshift=-2.8mm]ros.east);

    \node[auth] at (5.5, -1.45) {COMPUERTA DE AUTORIZACIÓN DE LABORATORIO Y SEGURIDAD: Requerida antes de cualquier ensayo en hardware[cite: 16]};
\end{tikzpicture}
\end{center}

\section{Matriz Comparativa de Rutas Candidatas A y B[cite: 16]}
Una ruta no es superior por popularidad, sino porque sus prerrequisitos documentados pueden comprobarse localmente con menor incertidumbre[cite: 16]:

\begin{center}
\footnotesize
\renewcommand{\arraystretch}{1.4}
\begin{tabularx}{\textwidth}{|>{\hsize=0.9\hsize\raggedright\arraybackslash}X|>{\hsize=1.05\hsize\raggedright\arraybackslash}X|>{\hsize=1.05\hsize\raggedright\arraybackslash}X|}
\hline
\rowcolor{ucbblue!15}
\textbf{Campo de Evaluación}[cite: 16] & \textbf{Ruta A (Integración con Driver Existente, ej. RWS/EGM condicionado)}[cite: 16] & \textbf{Ruta B (Puente Documentado Adaptado al Sistema Existente / RAPID)}[cite: 16] \\ \hline
1. Interfaz exacta candidata[cite: 16] & Por verificar[cite: 16] & Por verificar[cite: 16] \\ \hline
2. Versiones admitidas por fuente documental[cite: 16] & & \\ \hline
3. Prerrequisito local crítico (opción/licencia/código)[cite: 16] & & \\ \hline
4. Evidencia local actual del prerrequisito[cite: 16] & & \\ \hline
5. Canal de comando (conceptual)[cite: 16] & & \\ \hline
6. Canal de estado (origen, unidades, orden, timestamp)[cite: 16] & & \\ \hline
7. Prueba \textit{read-only} (solo lectura) propuesta[cite: 16] & & \\ \hline
8. Autorización de laboratorio necesaria[cite: 16] & & \\ \hline
9. Observable exacto y archivo de log[cite: 16] & & \\ \hline
10. Criterio \texttt{PASS / FAIL / BLOCKED}[cite: 16] & & \\ \hline
11. Condiciones de aborto y \texttt{NO-GO}[cite: 16] & & \\ \hline
12. Próxima acción y responsable asignado[cite: 16] & & \\ \hline
\end{tabularx}
\end{center}

\section{Especificación Estructurada del Ensayo Gate 1 (Diseñado, No Ejecutado)[cite: 16]}
\textbf{Pregunta verificable de Gate 1:} \textit{¿Puede la interfaz obtener, de manera autorizada y repetible, una señal real de estado o identidad del controlador ABB sin ordenar movimiento?}[cite: 16]

\begin{tcolorbox}[colback=white, colframe=ucbblue, title=\textbf{Protocolo de 8 Puntos para el Ensayo Gate 1} {[cite: 16]}, left=2mm, right=2mm, top=1mm, bottom=1mm]
    \footnotesize\raggedright
    \begin{enumerate}
        \item \textbf{Objetivo limitado:} Leer identificación o estado en modo \textbf{solo lectura}, sin envío de trayectoria ni activación de motores[cite: 16].
        \item \textbf{Precondiciones:} Protocolo/opción soportada, red autorizada, operador habilitado presente, modo de operación conforme al procedimiento local y separación estricta de interfaces de escritura[cite: 16].
        \item \textbf{Entrada:} Solicitud de consulta definida exactamente por la documentación de la versión concreta identificada (prohibido inventar llamadas de API)[cite: 16].
        \item \textbf{Observables requeridos:} Fuente verificable del estado, identificación del controlador, valor numérico y \textit{timestamp}, consistencia temporal, archivo de log, modo operativo y unidades[cite: 16].
        \item \textbf{Criterio \texttt{PASS}:} Respuesta rastreable originada en el controlador físico, reproducible y con ausencia total de movimiento. \textit{Alcance de la conclusión:} lectura de estado validada (no implica control de trayectoria)[cite: 16].
        \item \textbf{Criterio \texttt{FAIL}:} Intento autorizado ejecutado que no produce el observable esperado; exige conservar mensajes de error y configuración exacta[cite: 16].
        \item \textbf{Criterio \texttt{BLOCKED}:} Prerrequisito técnico o autorización de laboratorio no disponible; prohibido ejecutar la prueba o atribuir una falla técnica falsa[cite: 16].
        \item \textbf{Condición de Aborto:} Cualquier inconsistencia de identidad, error de permisos, comportamiento inesperado o instrucción directa del supervisor obliga a detener la prueba según el protocolo local sin improvisar \textit{bypasses}[cite: 16].
    \end{enumerate}
\end{tcolorbox}

\end{document}
